\documentclass{article}
\usepackage{amsmath,amssymb}
\usepackage[vocab={latticepoints}]{leanify}
\newtheorem{lemma}{Lemma}
\begin{document}

\section{Notation}

We write $N(r)$ for the number of lattice points in the disk of radius $r$.
We write $M(r)$ for the number of lattice points in the ball of radius $r$.
We write $S(a,b)$ for the number of lattice points on the segment from the origin to the point with coordinates $a$ and $b$.
We write $V(n)$ for the number of visible lattice points in the square of side $n$.
We write $A(i,b)$ for the area of a lattice polygon with $i$ interior points and $b$ boundary points.

\section{Counts}

\begin{lemma}\label{counts}
$N(1)=5$, $N(2)=13$, $M(1)=7$, $S(4,6)=3$ and $V(2)=3$.
\end{lemma}

\begin{lemma}\label{origin}
Let $r$ be a real number with $r\ge0$. Then $N(r)\ge1$.
\end{lemma}

\section{Pick's theorem}

\begin{lemma}\label{pick}
$A(0,4)=1$ and $A(1,4)=2$.
\end{lemma}

\begin{lemma}\label{pick-general}
Let $i$ and $b$ be integers with $i\ge0$ and $b\ge3$. Then $A(i,b)=i+\frac{b}{2}-1$.
\end{lemma}

\begin{lemma}\label{segment}
Let $a$ and $b$ be positive integers. Then $S(a,b)\ge2$.
\end{lemma}

\end{document}
